% Môn: Vật lý
% Lớp: 12
% Chương: Chương I. Vật lí nhiệt
% Bài: L12C1 Bài 4. Nhiệt dung riêng · Tính thời gian làm nóng vật
% ID: L12C1B3-121-TN
% ID: L12C1B3-121-TN
% Mức: NB
% ID: L12C1B3-121-TN


% Mức: VD
% Mức: VD
% ID: L12C1B3-136-TLN
% ID: L12C1B3-136-TLN
% Mức: VD
% ID: L12C1B3-136-TLN
% ID: L12C1B3-136-TLN
% Mức: VD
% ID: L12C1B4-77-TLN
%=== Câu 16 ===%
\dangbt{Tính thời gian làm nóng vật}
\begin{ex}
% ID: L12C1B4-77-TLN
% Mức: VD
Một vỉ đá có khối lượng $1,2\ kg$, nhiệt độ ban đầu là $28^\circ C$ được làm nóng trong lò có công suất $20\ kW$. Biết nhiệt dung riêng của đá là $5500\ J/kg.K$. Để vỉ đá đạt được nhiệt độ $1000^\circ C$ thì cần thời gian bao nhiêu phút (làm tròn kết quả đến chữ số hàng phần trăm)?
\shortans[oly]{$5,35$}
	\loigiai{
		\begin{itemize}
			\item $\mathcal{P} = \dfrac{Q}{t} \Rightarrow t = \dfrac{Q}{\mathcal{P}} = \dfrac{mc\Delta t}{\mathcal{P}} = \dfrac{1,2.5500.(1000-28)}{20.10^3} = 320,76\ s = 5,346$ phút.		
		\end{itemize}
	}
\end{ex}
